%% kapa: exam=1 round=24 year=2013 period=1 form=A subject=civil_law subject_ko=민법 seq=27 number=27
%% source: Q-Net 감정평가사 기출문제 (cst00309) artlSeq=5073802 file=extracted/1008/5073802/2013년 제24회 감정평가사 1차시험 기출문제/1교시(A형).hwp
%% conversion: kordoc 4.24.51 -> tools/gen1.py (auto); review=proofread-sonnet; flags=none
\begin{question}{27}
등기부시효취득에 관한 설명으로 옳지 않은 것은? (다툼이 있으면 판례에 의함)
\begin{choices}
  \item 무효인 이중보존등기에 기초해서는 등기부취득시효가 완성되지 않는다.
  \item 등기부시효취득에 있어서 무과실은 점유에 관한 것이고 등기에 관한 것이 아니다.
  \item 등기부시효취득에 있어서 과실 여부의 증명책임은 시효취득을 부인하는 소유자에게 있다.
  \item 지적공부 소관청의 분필절차를 거치지 않은 채 등기부상만으로 분할된 토지에 대한 등기부취득시효는 인정되지 않는다.
  \item 점유의 승계는 물론 등기의 승계도 인정한다.
\end{choices}

eviewflag{choices not rendered in rhwp page; only stem verified}
\end{question}
\answer{3}
